\clearpage
\ExposeTOC
\begin{center}
{\large SGA 7 --- 第 XX 讲}\par
\vspace{2em}
{\Large\underline{Griffiths 定理}}\par
\vspace{1.5em}
{\large N. Katz}
\end{center}

\section*{0. 引言}

我们证明 Griffiths 定理~[1], 它给出一种系统的 (但非构造性的) 方法, 用以``找到''在 $\mathbb Q_\ell$-上同调中等价于零、但其任何非零整数倍都不代数等价于零的代数圈. 这里给出的证明 (出自 Grothendieck) 是将 Griffiths 原来或多或少带有超越性质的证明转写为纯代数语言. 在这两个证明中, 不可约性定理 XVIII~6.7 均是关键 (参见~3.1.6 和~3.2).

\section*{1. 本原类的形式体系}

以下始终固定一个素数 $\ell$. 假定所考虑的所有概形的特征均不同于 $\ell$. 考虑下列情形

\begin{equation}\tag{1.0.1}
\begin{tikzcd}[column sep=huge, row sep=huge]
X \arrow[rr, "\pi"] \arrow[dr, "f"'] & {} & Y \arrow[dl, "g"] \\
{} & S & {}
\end{tikzcd}
\end{equation}
其中假设
\begin{equation}\tag{1.0.3}
f,\ g \quad\text{光滑, 相对维数分别为}\quad d_1\ \text{和}\ d_2,
\end{equation}
以及
\begin{equation}\tag{1.0.4}
\pi\quad\text{固有}.
\end{equation}
由于 $X$ 和 $Y$ 在 $S$ 上光滑, $\pi$ 的固有性给出一个\underline{Gysin 态射}(Gysin morphism) (SGA~5, 第~IV 讲; 印本中后续的引用留空)
\begin{equation}\tag{1.0.5}
\pi_*:H^q(X,\mathbb Q_\ell(i))\longrightarrow
H^{q-2(d_1-d_2)}(Y,\mathbb Q_\ell(i-d_1+d_2)).
\end{equation}
在 $S$ 的平展拓扑下局部化, 同样得到 $S$ 上层的态射\footnote{在 (1.0.6) 中, 印本以撇号代替指数中的右括号, 并将目标的扭写作 $i-d_1-d_2$. 公式 (1.0.5) 和 (1.0.7) 以及 Gysin 次数均给出此处所列的读法.}
\begin{equation}\tag{1.0.6}
\pi_*:R^qf_*\mathbb Q_\ell(i)\longrightarrow
R^{q-2(d_1-d_2)}g_*\mathbb Q_\ell(i-d_1+d_2),
\end{equation}
并在应用函子 $H^p(S,-)$ 后得到态射
\begin{equation}\tag{1.0.7}
\pi_*:H^p(S,R^qf_*\mathbb Q_\ell(i))\longrightarrow
H^p(S,R^{q-2(d_1-d_2)}g_*\mathbb Q_\ell(i-d_1+d_2)).
\end{equation}

\subsubsection*{定义 1.0.8.}
设 $(E_r^{p,q},d_r^{p,q})$ 和 $({}'E_r^{p,q},{}'d_r^{p,q})$ 为给定阿贝尔范畴中的两个谱序列, $(a,b)\in\mathbb Z\times\mathbb Z$. $({}'E_r^{p,q},{}'d_r^{p,q})$ 的 $(a,b)$ 移位是由下式给出的谱序列 $({}''E_r^{p,q},{}''d_r^{p,q})$\footnote{印本在 (1.0.9) 中漏去了右端微分上的撇号. 带撇号谱序列的移位的定义要求使用带撇号的微分.}
\begin{equation}\tag{1.0.9}
\left\{
\begin{aligned}
{}''E_r^{p,q}&={}'E_r^{p+a,q+b},\\
{}''d_r^{p,q}&={}'d_r^{p+a,q+b}.
\end{aligned}
\right.
\end{equation}

\subsubsection*{1.0.10.}
从 $(E_r^{p,q},d_r^{p,q})$ 到 $({}''E_r^{p,q},{}''d_r^{p,q})$ 的谱序列态射, 称为从 $(E_r^{p,q},d_r^{p,q})$ 到 $({}'E_r^{p,q},{}'d_r^{p,q})$ 的双次数为 $(a,b)$ 的态射.

下列引理是\underline{Gysin 态射}理论 (SGA~5, 第~IV 讲) 的直接推论.

\begin{editorialnote}
译校: 引理中的 $\alpha:S\to T$ 是下文图 (1.2.1) 中的态射; 法文印本第 343 页在引理之前尚未引入此记号. 公式 (1.0.5)、(1.0.7) 给出绝对上同调的写法; 引理中的相对形式则相应地用 $R\alpha_*$ 及 $R(\alpha f)_*$、$R(\alpha g)_*$ 表示, 如引理的两个 Leray 谱序列所示.
\end{editorialnote}

\subsubsection*{引理 1.1.}
在 Leray 谱序列
\begin{equation}\tag{1.1.0}
R^p\alpha_*R^qf_*\mathbb Q_\ell(i)\Longrightarrow
R^{p+q}(\alpha f)_*\mathbb Q_\ell(i)
\end{equation}
和
\begin{equation}\tag{1.1.1}
R^p\alpha_*R^qg_*\mathbb Q_\ell(i-d_1+d_2)\Longrightarrow
R^{p+q}(\alpha g)_*\mathbb Q_\ell(i-d_1+d_2),
\end{equation}
之间, 存在双次数为 $(0,-2(d_1-d_2))$ 的典范态射 $\pi_*$\footnote{印本写作 $(a,-2(d_1-d_2))$. 映射 (1.0.7) 不改变第一个 Leray 指标 $p$, 因而双次数的第一个分量为 $0$.}, 它在 $E_2$ 项上由 (1.0.7) 给出, 在极限上由 (1.0.5) 给出.

\subsubsection*{1.2.}
现在考虑下列情形
\begin{equation}\tag{1.2.1}
\begin{tikzcd}[column sep=large, row sep=large]
{} & X\times_S Y \arrow[dl, "\mathrm{pr}_1"']
                  \arrow[dr, "\mathrm{pr}_2"] & {} \\
X \arrow[dr, "f"'] & {} & Y \arrow[dl, "g"] \\
{} & S \arrow[d, "\alpha"] & {} \\
{} & T & {}
\end{tikzcd}
\end{equation}
其中 $S,X,f,Y,g$ 如 (1.0.1)--(1.0.4), 但现在还假设 $f$ 固有.\footnote{印本将此范围的起点写作不存在的 (1.0.2); 所列的情形为 (1.0.1).} 设
\begin{equation}\tag{1.2.2}
z\in H^a(X\times_S Y,\mathbb Q_\ell(b)).
\end{equation}
则 $z$ 诱导仍记为 $z$ 的映射
\begin{equation}\tag{1.2.3}
z:H^q(X,\mathbb Q_\ell(i))\longrightarrow
H^{q+a-2d_1}(Y,\mathbb Q_\ell(i+b-d_1))
\end{equation}
其公式为
\begin{equation}\tag{1.2.4}
z(x)=\mathrm{pr}_{2*}\bigl(z\wedge\mathrm{pr}_1^*(x)\bigr),
\end{equation}
其中 $\mathrm{pr}_{2*}$ 为 Gysin 同态. 在 $S$ 上作平展局部化, 类 $z$ 同样定义
\begin{equation}\tag{1.2.5}
z:R^qf_*\mathbb Q_\ell(i)\longrightarrow
R^{q+a-2d_1}g_*\mathbb Q_\ell(i+b-d_1),
\end{equation}
并在应用 $H^p(S,-)$ 后定义映射
\begin{equation}\tag{1.2.6}
H^p(S,R^qf_*\mathbb Q_\ell(i))\longrightarrow
H^p(S,R^{q+a-2d_1}g_*\mathbb Q_\ell(i+b-d_1)).
\end{equation}

\subsubsection*{引理 1.3.}
在 Leray 谱序列
\begin{equation}\tag{1.2.7}
H^p(S,R^qf_*\mathbb Q_\ell(i))\Longrightarrow
H^{p+q}(X,\mathbb Q_\ell(i))
\end{equation}
和
\begin{equation}\tag{1.2.8}
H^p(S,R^qg_*\mathbb Q_\ell(i+b-d_1))\Longrightarrow
H^{p+q}(Y,\mathbb Q_\ell(i+b-d_1)),
\end{equation}
之间, 类 $z$ 定义双次数为 $(0,a-2d_1)$ 的态射\footnote{印本写作 $(0,b-2d_1)$. 映射 (1.2.3), (1.2.5) 和 (1.2.6) 将上同调次数改变 $a-2d_1$; $b$ 控制 Tate 扭.}, 它在 $E_2$ 项上由 (1.2.6) 给出, 在极限上由 (1.2.3) 给出.

\noindent\underline{证明.}
由 (1.2.4), 又因``乘以 $z$''在 Leray 谱序列
\begin{equation}\tag{1.2.9}
H^p\bigl(S,R^q(f\circ\mathrm{pr}_1)_*\mathbb Q_\ell(i)\bigr)
\Longrightarrow H^{p+q}(X\times_S Y,\mathbb Q_\ell(i))
\end{equation}
\begin{center}
\begin{tikzcd}
{} \arrow[d, "z"'] \\
{}
\end{tikzcd}
\end{center}
\begin{equation}\tag{1.2.10}
H^p\bigl(S,R^{q+a}(f\circ\mathrm{pr}_1)_*\mathbb Q_\ell(i+b)\bigr)
\Longrightarrow H^{p+q+a}(X\times_S Y,\mathbb Q_\ell(i+b)),
\end{equation}
之间定义双次数为 $(0,a)$ 的态射, 故只须对下列情形应用引理~1.1\footnote{在 (1.2.9) 和 (1.2.10) 中, 印本写作 $(\mathrm{pr}_1 f)_*$, 但到 $S$ 的映射为 $f\circ\mathrm{pr}_1$, 正如印本的图 (1.2.11) 本身所示.}
\begin{equation}\tag{1.2.11}
\begin{tikzcd}[column sep=large, row sep=large]
X\times_S Y \arrow[rr, "\mathrm{pr}_2"]
             \arrow[dr, "f\circ\mathrm{pr}_1"]
  & {} & Y \arrow[dl, "g"] \\
{} & S & {}
\end{tikzcd}
\qquad .
\end{equation}

\subsubsection*{推论 1.4.}
在 (1.2.1)--(1.2.2) 的情形下, 态射 (1.2.3)
\begin{equation}\tag{1.4.1}
z:H^q(X,\mathbb Q_\ell(i))\longrightarrow
H^{q+a-2d_1}(Y,\mathbb Q_\ell(i+b-d_1))
\end{equation}
将本原类 (XVIII~5.8.1) 映为本原类, 且下图交换:
\begin{equation}\tag{1.4.2}
\begin{tikzcd}[column sep=large, row sep=large]
\mathrm{Prim}^q(X/S,\mathbb Q_\ell(i))
  \arrow[r, "z"] \arrow[d, "u_f"]
& \mathrm{Prim}^{q+a-2d_1}(Y/S,\mathbb Q_\ell(i+b-d_1))
  \arrow[d, "u_g"] \\
H^1(S,R^{q-1}f_*\mathbb Q_\ell(i))
  \arrow[r, "z"]
& H^1(S,R^{q+a-2d_1-1}g_*\mathbb Q_\ell(i+b-d_1)).
\end{tikzcd}
\end{equation}
这里 $u_f$ 和 $u_g$ 是 XVIII~5.8.1 中的典范同态.
